\section{全概率公式与 Bayes 更新}
\label{sec:05-Probability/02-Conditional-Probability/02}

你管理一家工厂，两台机器日夜不停地生产同一种零件。机器 A 产量占 \(70\%\)，技术成熟，次品率只有 \(1\%\)。机器 B 产量占 \(30\%\)，设备老化，次品率 \(4\%\)。所有的零件混在一起进入仓库，你随手拿起一个——是次品。

问题来了：这只次品更可能是哪台机器生产的？

你可能会想：机器 A 产量大，基数大，次品更可能来自 A。但转念一想：机器 B 的次品率是 A 的四倍，虽然产量小，每生产一百个就出四个坏零件。两个力量在拉扯——基数 vs 次品率。谁赢？

这就是 Bayes 更新的核心场景：你有一个看不见的原因（机器 A 还是 B），你只能观察到它产生的结果（零件是好是坏）。你知道从原因到结果的概率——\(P(\text{次品}\mid A)=0.01\)，\(P(\text{次品}\mid B)=0.04\)。但你真正想知道的是从结果反推原因：\(P(A\mid\text{次品})\) 和 \(P(B\mid\text{次品})\)。条件概率的竖线不能直接交换，但 Bayes 公式给你一座桥。

\subsection{全概率公式：先搞清楚次品总共有多常见}

在反推之前，先算一个边际量：随机抽一个零件，它是次品的总概率有多大？

次品只能沿两条互不重叠的路径产生：要么由机器 A 生产且通不过质检，要么由机器 B 生产且通不过质检。两条路径穷尽了所有可能：A 和 B 覆盖全部产量。

沿第一条路径：由 A 生产且为次品。\(P(A)=0.7\)，在 A 生产的条件下次品的概率 \(P(D\mid A)=0.01\)。所以该路径的概率为 \(0.7\times0.01=0.007\)。

沿第二条路径：由 B 生产且为次品。\(P(B)=0.3\)，\(P(D\mid B)=0.04\)。概率为 \(0.3\times0.04=0.012\)。

总次品率：

\[
P(D)=0.7\times0.01+0.3\times0.04=0.007+0.012=0.019.
\]

约 \(1.9\%\)。注意你不能把 \(1\%\) 和 \(4\%\) 简单平均为 \(2.5\%\)——那样忽略了产量权重。机器 A 的低次品率有更大的权重（\(70\%\)），所以总次品率被拉低到了 \(1.9\%\)，比简单平均 \(2.5\%\) 更接近 A 的次品率。

把这个推理一般化：若事件 \(H_1,\ldots,H_m\) 两两互斥并覆盖样本空间（即 \(\bigcup_i H_i=\Omega\)），任意事件 \(E\) 的概率可以拆成互斥路径之和：

\[
P(E)=\sum_{i=1}^{m}P(H_i)P(E\mid H_i).
\]

这就是\textbf{全概率公式}。它做的事可以用一句话描述：\textbf{把结果 \(E\) 的所有来路找出来，每条路用``走到这条路的概率''乘上``在这条路上发生 \(E\) 的概率''，然后加总}。

你可能已经注意到：这不就是链式法则和互斥事件可加性的组合吗？确实是。\(P(H_i)P(E\mid H_i)=P(E\cap H_i)\)，而 \(E=\bigcup_i(E\cap H_i)\) 是互斥并。全概率公式没有引入新公理，只是把一个事件按它的``原因''做了划分。

\textbf{划分不对，加总白费。}全概率公式要求 \(H_1,\ldots,H_m\) 是一个划分——两两互斥且并集为全集。如果遗漏了一个原因（比如工厂还有外包机器 C 你没算），分母 \(P(E)\) 就算小了，后续的 Bayes 后验也跟着偏移。如果在划分中把同一个结果重复分配到了两条路径（即 \(H_i\) 有重叠），概率会被多算。检查划分是否穷尽且互斥，是任何全概率计算之前必须做的第一件事。

\subsection{Bayes 公式：反推原因}

现在可以回答最初的问题了。已知零件是次品，它来自机器 B 的概率是多少？

由条件概率定义，

\[
P(B\mid D)=\frac{P(B\cap D)}{P(D)}.
\]

分子 \(P(B\cap D)=0.3\times0.04=0.012\)。分母就是刚才全概率公式算出的 \(P(D)=0.019\)。所以

\[
P(B\mid D)=\frac{0.012}{0.019}\approx0.632.
\]

机器 B 只生产了 \(30\%\) 的零件，但贡献了约 \(63.2\%\) 的次品。基数被次品率翻盘了。

类似地，来自 A 的概率：

\[
P(A\mid D)=\frac{0.7\times0.01}{0.019}=\frac{0.007}{0.019}\approx0.368.
\]

两个后验概率加起来是 \(1\)——次品总得来自某台机器。原来是 \(0.7\) vs \(0.3\)（先验），现在是 \(0.368\) vs \(0.632\)（后验）。证据``这是个次品''把概率质量从 A 推向 B。方向对了，因为 B 的次品率更高。

你可以把这个过程想象成一场拔河。初始时，绳子上的标记偏向 A——\(70\%\) 的产量，A 占据绝对优势。然后证据出现了：这只零件是次品。证据在 B 的方向上拉了一把——因为 B 的次品率是 A 的四倍，证据对 B 更``友好''。拉完之后，标记移到了 \(63.2\%\) 偏向 B 的位置。拉绳的力度取决于似然比 \(P(D\mid B)/P(D\mid A)=0.04/0.01=4\)。

把上述计算写成一般形式，就是\textbf{Bayes 公式}：

\[
P(H_i\mid E)=\frac{P(H_i)P(E\mid H_i)}{\sum_j P(H_j)P(E\mid H_j)}.
\]

可以把它读成一个三拍子的节奏：

\[
\text{后验}\propto\text{先验}\times\text{似然}.
\]

\begin{itemize}
\tightlist
\item \textbf{先验} \(P(H_i)\)：在你看到任何证据之前，你对原因 \(H_i\) 赋予的权重。工厂例子中是 \(0.7\) 和 \(0.3\)。
\item \textbf{似然} \(P(E\mid H_i)\)：如果原因 \(H_i\) 是真的，证据 \(E\) 出现的自然程度。工厂例子中是 \(0.01\) 和 \(0.04\)。
\item \textbf{后验} \(P(H_i\mid E)\)：在看到了证据 \(E\) 之后，你对原因 \(H_i\) 重新分配的权重。分母 \(\sum_j P(H_j)P(E\mid H_j)\) 即 \(P(E)\)，负责把后验概率重新规范化为总和 \(1\)。
\end{itemize}

你可能会想：这和日常思考有什么不同？日常思考也在用新信息更新判断。但日常思考不会显式地标出先验从哪来、似然基于什么数据、分母有没有穷尽所有可能原因。Bayes 公式把更新过程变成了透明的运算，每一步都可以被检查、被质疑、被修正。

\subsection{为什么阳性之后你不一定该紧张}

一种疾病在人群中的患病率是 \(1\%\)。检测方法的技术指标如下：对真正患病的人，有 \(99\%\) 的概率测出阳性（敏感性）；对健康的人，有 \(5\%\) 的概率也测出阳性（假阳性率）。

你去做了一次检测，结果是阳性。你有多大可能真的患病？

你的第一反应可能是``\(99\%\)？那基本就是了''。你可能会这么想：检测对病人有 \(99\%\) 的正确率，那我阳性的情况下，\(99\%\) 就是病人。但等一下——你混淆了方向。\(99\%\) 是 \(P(+\mid D)\)（病人中测出阳性的概率），而你想知道的是 \(P(D\mid+)\)（阳性者中真正患病的概率）。这是两个不同的问题，竖线两边的角色完全相反。

让我们用刚才的三拍子节奏走一遍。

两个原因：患病（\(D\)）和健康（\(H\)）。先验：\(P(D)=0.01\)，\(P(H)=0.99\)。似然：\(P(+\mid D)=0.99\)，\(P(+\mid H)=0.05\)。

全概率分母——随机一个人检测出阳性的总概率：

\[
P(+)=0.01\times0.99+0.99\times0.05=0.0099+0.0495=0.0594.
\]

注意：虽然健康人的假阳性率只有 \(5\%\)，但健康人占了人群的 \(99\%\)，所以来自健康人的阳性贡献 \(0.0495\) 远远超过了来自病人的 \(0.0099\)。假阳性的绝对数量淹没了真阳性。

Bayes 后验：

\[
P(D\mid+)=\frac{0.01\times0.99}{0.0594}=\frac{0.0099}{0.0594}\approx0.167.
\]

只有约 \(16.7\%\)。检测阳性，但真正患病的概率不到两成。

为什么会这么低？不是检测不准——\(99\%\) 的敏感性确实很高——而是健康人的基数太大了。一万个人里，约 \(100\) 人患病，其中 \(99\) 人阳性；约 \(9900\) 人健康，其中 \(5\%\) 也就是 \(495\) 人假阳性。所以总共约 \(594\) 人拿到阳性报告，而其中真正患病的只有 \(99\) 人。\(99/594\approx16.7\%\)。

把抽象概率翻译成自然人数——一万个人里的实际人头数——许多人的困惑瞬间消失。你不需要更高级的公式，只需要把``\(1\%\)''变成``一百人''，把``\(99\%\)''变成``九十九人''，把``\(5\%\)''变成``四百九十五人''。往后每次被 Bayes 公式绕晕时，试试画一个 \(10,000\) 人的表格。

但也要警惕：这个 \(16.7\%\) 极度依赖于 \(1\%\) 的患病率。如果患病率是 \(10\%\) 而不是 \(1\%\)——比如高风险人群——那么

\[
P(D\mid+)=\frac{0.1\times0.99}{0.1\times0.99+0.9\times0.05}=\frac{0.099}{0.144}\approx0.688.
\]

同一个检测，在普通人群中阳性后患病概率 \(16.7\%\)，在高风险人群中 \(68.8\%\)。先验变了，后验跟着大翻盘。这就是为什么医生在开检测单之前要问家族史和行为风险——先验不是随便填的数字。

换个极端方向再感受一下：如果假阳性率从 \(5\%\) 改善到 \(0.1\%\)（一个极其精确的检测），普通人群中 \(P(D\mid+)=(0.01\times0.99)/(0.01\times0.99+0.99\times0.001)\approx0.0099/0.01089\approx0.909\)。同样的疾病患病率，同样的敏感性，假阳性率从 \(5\%\) 降到 \(0.1\%\) 后，阳性后患病概率从 \(16.7\%\) 跳到了 \(90.9\%\)。检测的假阳性率对后验的影响，在某些先验下甚至比敏感性更关键。

\subsection{经典场景：Monty Hall 问题}

你参加一个电视游戏节目，面前有三扇关闭的门。其中一扇后面是一辆车，另外两扇后面各是一只山羊。你选一扇门（比如 \(1\) 号），但先不打开。主持人 Monty——他知道每扇门后面是什么——会从你未选的两扇门中，打开一扇后面是山羊的门（比如 \(3\) 号）。然后他问你：你要坚持原来的选择（\(1\) 号），还是换到剩下的那扇门（\(2\) 号）？

你可能会想：剩下两扇门，一扇有车一扇有山羊，换不换概率都是 \(50\%\)，没区别。这个直觉非常强大，强大到很多数学家在第一次听到这个问题时都犯了错。

但如果用 Bayes 更新来做呢？

设 \(C_i\) 表示车在 \(i\) 号门后。先验：\(P(C_1)=P(C_2)=P(C_3)=1/3\)。你初选了 \(1\) 号门。Monty 打开了 \(3\) 号门，放出一只山羊。我们想知道 \(P(C_1\mid \text{Monty开}3)\) 和 \(P(C_2\mid \text{Monty开}3)\)。

关键在一个容易被忽略的细节：Monty 不是随机开门的。如果车在 \(1\) 号门后（你的初选），Monty 可以在 \(2\) 和 \(3\) 中任选一扇——假设他各选一半。如果车在 \(2\) 号门后，Monty 只能开 \(3\) 号（因为 \(1\) 是你选的不能开，\(2\) 有车不能开）。如果车在 \(3\) 号门后，Monty 只能开 \(2\) 号——但他开了 \(3\) 号，所以这种情况不可能。

因此：
\[
\begin{aligned}
P(\text{Monty开}3\mid C_1)&=\frac12,\\
P(\text{Monty开}3\mid C_2)&=1,\\
P(\text{Monty开}3\mid C_3)&=0.
\end{aligned}
\]

全概率分母：
\[
P(\text{Monty开}3)=\frac13\cdot\frac12+\frac13\cdot1+\frac13\cdot0=\frac16+\frac13=\frac12.
\]

Bayes 后验：
\[
\begin{aligned}
P(C_1\mid\text{Monty开}3)&=\frac{(1/3)\times(1/2)}{1/2}=\frac13,\\[4pt]
P(C_2\mid\text{Monty开}3)&=\frac{(1/3)\times1}{1/2}=\frac23.
\end{aligned}
\]

换门后赢车的概率是 \(2/3\)，坚持原选只有 \(1/3\)。差距不在车的位置变了——车始终在最初那扇门后面——而在于 Monty 的开门行为携带着信息：他避开了你的初选门，也避开了有车的门。当你的初选是错误的时候，他只有唯一一扇门可开；当你的初选正确时，他有选择余地。这个不对称性通过似然 \(P(\text{Monty开}3\mid C_i)\) 进入了后验，把概率推向了换门的一方。

Monty Hall 教会你两件事：(1) Bayes 公式中，似然常常比你直觉中认为的更复杂——它不仅取决于结果，还取决于产生结果的机制和约束。(2) 当你的答案和直觉严重冲突时，画出信息流——谁在什么时候知道了什么——往往比直接算概率更快找到症结。

\subsection{赔率形式：乘上似然比，约掉分母}

当只有两个假设 \(H_1\) 和 \(H_0\) 时，把各自的 Bayes 后验相除：

\[
\frac{P(H_1\mid E)}{P(H_0\mid E)}
=\frac{P(H_1)}{P(H_0)}\cdot\frac{P(E\mid H_1)}{P(E\mid H_0)}.
\]

分母 \(P(E)\) 在相除时约掉了，省掉一步归一化。

左端是\textbf{后验赔率}（posterior odds）：看到证据之后，\(H_1\) 相对于 \(H_0\) 的可能性倍数。右端第一个因子是\textbf{先验赔率}（prior odds）；第二个因子是\textbf{似然比}（likelihood ratio），衡量证据在两种假设下出现的可能性之比。

赔率形式的好处在你连续收到多条证据时尤其明显。如果两条证据 \(E_1,E_2\) 在给定假设下条件独立（条件独立的精确定义在本单元\hyperref[sec:05-Probability/02-Conditional-Probability/04]{``条件独立与信息结构''}中给出），那么：

\[
\frac{P(H_1\mid E_1,E_2)}{P(H_0\mid E_1,E_2)}
=\frac{P(H_1)}{P(H_0)}\cdot\frac{P(E_1\mid H_1)}{P(E_1\mid H_0)}\cdot\frac{P(E_2\mid H_1)}{P(E_2\mid H_0)}.
\]

每一次新证据只是在上一次的赔率上再乘一个似然比。取对数之后，证据的贡献直接相加——每一条新信息用一个加数更新你的对数赔率。这就是许多分类器、序贯检验和 Bayesian 统计推断的核心运算结构。

但有一个陷阱：``重复检测''只有在给定真实状态后近似条件独立时，似然比才能直接相乘。如果两次检测用的是同一台有系统偏差的仪器，它们的错误方向可能相关——第一次假阳性意味着第二次假阳性的概率可能也高于 \(5\%\)。此时它们并没有提供两份独立证据，赔率乘积公式不再成立。

用数字说明：如果两次检测完全条件独立给定患病状态，且指标都是敏感性 \(99\%\)、假阳性 \(5\%\)，那么初始患病率 \(1\%\) 的情况下，一次阳性后赔率从 \((0.01/0.99)\approx0.0101\) 变为 \(0.0101\times(0.99/0.05)\approx0.2\)（即后验概率约 \(16.7\%\)），第二次阳性后再乘一次 \(0.99/0.05=19.8\)，赔率变为 \(0.2\times19.8\approx3.96\)，后验概率约 \(3.96/4.96\approx79.8\%\)。两次阳性后你开始真正担心了。但如果两次检测共享系统偏差，第二次乘的因子可能远小于 \(19.8\)，你实际的后验概率远低于 \(79.8\%\)。

\subsection{Bayes 公式不替你检查模型}

Bayes 更新可以把一堆输入数字变成输出数字，但你喂进去的数字本身是否靠得住，公式不管。以下每一项出问题，都会让后验变成精确的错误：

\begin{itemize}
\tightlist
\item \textbf{候选原因是否穷尽。}工厂例子中，如果你的零件还可能是外包厂 C 生产的，但你只放了 A 和 B 两个假设，后验概率会在两个错误选项之间分配，给不出正确答案。
\item \textbf{先验从哪来。}\(P(D)=0.01\) 是基于全国普查、某个医院的历年数据，还是你拍脑袋填的？先验的微弱差异可能被高似然比放大成后验的巨大偏移。
\item \textbf{似然的适用人群。}检测的 \(99\%\) 敏感性和 \(5\%\) 假阳性率是在什么人群中测量的？如果在年轻人中测的指标，用在老年人身上，你的似然可能就不对。
\item \textbf{证据间的依赖结构。}如上一小节所说，把相关证据当成独立证据，后验看起来会比你实际该有的信心更极端。
\end{itemize}

所以一个负责任的 Bayes 更新，不仅要给出后验数字，还要标明数字的前提条件。当有人给你一个后验概率时，你的第一反应应该是问：先验是怎么定的？似然数据来源于哪里？还有没有其他可能的解释？这份怀疑不是对 Bayes 公式的不敬——恰恰是 Bayes 推理的组成部分。

\subsection{你从这里出发}

全概率公式和 Bayes 公式构成了一对工具：一个顺着原因到结果的方向（全概率），一个反着从结果回到原因（Bayes）。它们绑在一起的纽带是条件概率的定义——两者都不需要新公理。

它们不回答``为什么''——不解释因果关系，不判断哪个假设更好——只回答一个运算问题：在给定的先验和似然下，证据如何改变你对不同可能原因的概率分配。

但如果你走到了这一步，接下来该问的问题是：很多概率公式在做乘法时，其实默认了一个假设——事件之间互不影响。这个假设什么时候成立？什么时候不成立？怎么用公式精确地表达``互不影响''？

下一节讨论独立性。独立性是一个比``不相关''、``没关系''强得多的概念：它要求联合概率精确地分解为边缘概率的乘积。许多你在概率论里会反复见到的结论——大数定律、中心极限定理、随机过程的简化——都把它当作支点。而在独立性不成立的方向上，条件独立为 Bayes 网络的骨架、图模型的结构学习、以及高维数据的稀疏建模铺开了道路。
